%% ma: L12-B08-0178
%% lop: 12
%% chuong: 2
%% bai: 8
%% dang: 12.8.53
%% loai: TN4
%% muc: VD
%% dap_an: "A"
%% nguon: "CĐ2.4.pdf giáo viên gửi 27/10/2026 (M07, Lê Minh Khôi – Nguyễn Anh Khôi, Chương II Chủ đề 4 Các bài toán về tâm tỉ cự), Câu 28"
%% kien_thuc: "Bài 8 (khoảng cách hai điểm, biểu thức tọa độ), Bài 6 (bình phương vô hướng, tâm tỉ cự), Bài 7 (hình chiếu lên mặt phẳng tọa độ)"
%% trang_thai: cho-duyet
%% ly_do: "Dạng mới 12.8.53 đề xuất, chờ giáo viên duyệt"
\begin{ex}
Trong không gian $Oxyz$, cho ba điểm $A(2;0;-1)$, $B(1;0;-1)$, $C(0;1;0)$. Gọi $M$ là điểm thuộc mặt phẳng $(Oxy)$ sao cho $AM^2-5BM^2+2CM^2$ đạt giá trị lớn nhất. Tính độ dài đoạn thẳng $OM$.
\choice{\True $OM=\dfrac{\sqrt{13}}2$}{$OM=\dfrac{\sqrt{29}}2$}{$OM=\dfrac{\sqrt{26}}2$}{$OM=\dfrac{\sqrt6}2$}
\loigiai{Gọi $I$ thỏa $\overrightarrow{IA}-5\overrightarrow{IB}+2\overrightarrow{IC}=\vec 0\Rightarrow I=\dfrac{A-5B+2C}{-2}=\left(\dfrac32;-1;-2\right)$.

Biểu thức bằng $-2MI^2+$ hằng số, lớn nhất khi $M\left(\dfrac32;-1;0\right)$. Vậy $OM=\dfrac{\sqrt{13}}2$.}
\end{ex}
